\documentclass[10pt,a4paper]{amsart}
\usepackage[margin=19mm]{geometry}
\usepackage{amsmath,amssymb,mathtools}
\usepackage[scr=rsfs,cal=euler]{mathalfa}
\usepackage{xfrac}
\usepackage[unicode,hidelinks,bookmarks=false]{hyperref}
\newcommand{\ie}{i.e.\ }
\newcommand{\cf}{cf.\ }
\newcommand{\resp}{respectively\ }
\newcommand{\Subsection}{\subsection}
\providecommand{\nobreakdash}{\nobreak\hbox{-}}
\setlength{\emergencystretch}{2em}
\multlinegap0pt
\usepackage{mathtools}
\usepackage{amssymb}
\usepackage[shortlabels]{enumitem}
\usepackage{tikz}
\newtheorem{proposition}{Proposition}
\newcommand{\EE}{\mathbb{E}}
\newcommand{\N}{\mathbb{N}}
\newcommand{\PP}{\mathbb{P}}
\newcommand{\X}{\mathcal{X}}
\title{Relation between Hitting Times and Probabilities for Imprecise Markov Chains}
\author{Marco Sangalli, Erik Quaeghebeur, Thomas Krak}
\date{Source: arXiv:2603.16675v1 (CC BY 4.0)}
\begin{document}
\maketitle

\noindent\emph{Source and licence.} Relation between Hitting Times and Probabilities for Imprecise Markov Chains. Version arXiv:2603.16675v1 (CC BY 4.0), distributed by arXiv under the Creative Commons Attribution 4.0 International licence, http://creativecommons.org/licenses/by/4.0/, as recorded in the arXiv licence metadata for this exact version and shown on the abstract page https://arxiv.org/abs/2603.16675v1. Full licence notice: \texttt{licenses/arxiv-2603.16675-CC-BY-4.0.txt}.\par\smallskip

\noindent\textit{Editorial scope.} Counted unit: the article's proposition prop:leftimpl (source0.tex lines 190--195). The article's complete original proof of it is delivered with it (source0.tex lines 196--223, one \texttt{proof} environment of 28 lines). No mathematics is written, repaired or re-derived: the statement and the proof are the article's own lines, in the article's own notation, and no retained line is substituted or re-flowed.  No further source-local context is needed: every cross-reference of every retained line resolves inside the two delivered spans.  Nothing was omitted from any retained span, so the package carries no omission record.  No retained line contains a citation, so the wrapper carries no bibliography and no bibliography entry.  No retained line contains a figure environment, an image-inclusion command or drawing code, so no image is delivered and the package declares no asset dependency.  Editorial formatting only, in addition: an AMS \texttt{amsart} preamble that restates the article's own declarations and macros used by the retained lines, namely the environments \texttt{proposition} on separate counters, together with the standard packages those lines need.  The retained environments print in this document as Proposition 1 in order of appearance, whereas the article prints the same items with the numbering of its own full document.  The complete native source of the article as distributed in the arXiv e-print for this exact version is bundled unchanged as \texttt{sources/196568/source0.tex}.
\begin{proposition}\label{prop:leftimpl}
    Let $\X$ be a finite state space and $T$ be a transition matrix. Then 
    \begin{equation}
        h^T(x)<+\infty \ \Leftarrow \ p^T(x)=1.
    \end{equation}
\end{proposition}

\begin{proof}
    Let $R(x)\coloneqq x \cup \{y\in \X: \exists n\in \N, \ \PP_T(X_n=y \mid X_0=x)>0\}.$ Fix $n\in \N$. Then
    \begin{align}
        1&=p^T(x)=\PP_T(\tau_A<+\infty\mid X_0=x)\nonumber\\
        &= \EE_{\PP_T}\left[\PP_T(\tau_A<+\infty \mid X_n) \mid X_0=x \right] \nonumber\\
        &=\sum_{y\in \X}\PP_T(X_n=y\mid X_0=x)\PP_T(\tau_A<+\infty\mid X_0=y), \label{rompicapo}
    \end{align}
    where we used the law of total expectation and Markov's property. It follows that, for all $y\in \X$ with $\PP_T(X_n=y\mid X_0=x)>0$ it must hold that $\PP_T(\tau_A<+\infty\mid X_0=y)=p^T(y)=1$. For the arbitrariness of $n$, we get $p^T(y)=1$ for all $y\in R(x)$.

    Fix $y\in R(x)$. Since $p^T(y)= \PP_T(\tau_A<+\infty\mid X_0=y)=1$, for all $\varepsilon>0$, there exists $m_{y,\varepsilon}\in\N$ such that for all $n\ge m_{y,\varepsilon}$ we have $\PP_T(\tau_A\le n\mid X_0=y)\ge 1-\varepsilon$ or, equivalently, $\PP_T(\tau_A> n\mid X_0=y)\le\varepsilon$.
Fix $0<\varepsilon<1$ and set $M=\max_{y\in R(x)} m_{y,\varepsilon}$. Then, by picking $n=M$, we have
\[
\PP_T(\tau_A> M\mid X_0=y)\le\varepsilon,
\]
for all $y\in R(x)$.
Let $k\in \N_0$. Then, using Markov's property, we have
\[
\PP_T(\tau_A>kM\mid X_0=x)\le\varepsilon^k.
\]
We can bound the expected value of $\tau_A$ using these tail bounds:
\begin{align*}
    h^T(x)&=\EE_{\PP_T}[\tau_A\mid X_0=x] = {\sum_{n\in \N_0}\PP_T(\tau_A>n\mid X_0=x)}\\
&= \sum_{k\in \N_0}\sum_{n=kM}^{(k+1)M-1}\PP_T(\tau_A>n\mid X_0=x)\\
&\le \sum_{k\in \N_0} M\,\PP_T(\tau_A>kM\mid X_0=x)\\
&\le M\sum_{k\in \N_0}\varepsilon^k= \dfrac{M}{1-\varepsilon}<+\infty. 
\end{align*}
\qed
\end{proof}

\end{document}
